\documentclass[../../../include/open-logic-section]{subfiles}

\begin{document}

\olfileid[pt]{his}{set}{hilbertcurve}

\olsection{Apêndice: Curvas Preenchedoras de Espaço de Hilbert}

No capítulo~\olref[pathology]{sec}, mencionamos que a prova de Cantor de que
uma reta e um quadrado têm exatamente o mesmo número de pontos
(\olref[his][set][cantorplane]{thm:cantorplane}) levou Peano a
perguntar se poderia existir uma \emph{curva} que preenchesse o espaço. Obteve
resposta afirmativa em 1890. Nesta secção explicamos (de modo informal) como
construir a curva preenchedora de espaço de Hilbert (com um pequeno
ajuste).\footnote{Para uma explicação mais rigorosa, ver
\cite{Rose2010}. O ajuste consiste em incluir as partes vermelhas
das curvas abaixo. Isto torna ligeiramente mais fácil verificar que
a curva é contínua.}

Temos de definir uma função, $h$, como limite de uma sequência de funções $h_1$, $h_2$, $h_3$, \dots\@ Primeiro descrevemos a construção. Depois mostramos que preenche o espaço. Depois mostramos que é uma curva.

Tomaremos a imagem de~$h$ como sendo o quadrado unitário, $\unitsquare$. Eis a nossa primeira aproximação a~$h$, O.s.\@, $h_1$:
\begin{center}
	\begin{tikzpicture}
	\draw[gray] (0pt, 0pt) rectangle (80pt, 80pt);
	\draw[step = 40pt, gray, very thin] (0pt, 0pt) grid (80pt, 80pt);
	\draw[oldiagcolorC, thick] (20pt, 0)--(20pt, 20pt);
	\draw[oldiagcolorC, thick] (60pt, 0)--(60pt, 20pt);		
	\begin{scope}[xshift=20pt, yshift=20pt]
	\draw [black, thick, l-system={Hilbert curve, step=40pt, angle=90, axiom=L, order=1}]  lindenmayer system;
	\end{scope}
	\end{tikzpicture}
\end{center}
Para orientação, impusemos uma grelha $2 \times 2$ ao quadrado. Podemos pensar na curva como começando no quadrante inferior esquerdo, indo para o superior esquerdo, depois para o superior direito e finalmente para o inferior direito. Eis a segunda etapa da construção, O.s.\@, $h_2$:
\begin{center}
	\begin{tikzpicture}
	\draw[gray] (0pt, 0pt) rectangle (80pt, 80pt);
	\draw[step = 20pt, gray, very thin] (0pt, 0pt) grid (80pt, 80pt);
	\draw[oldiagcolorC, thick] (0pt, 10pt)--(10pt, 10pt);
	\draw[oldiagcolorC, thick] (70pt, 10pt)--(80pt, 10pt);
	\draw[thick, oldiagcolorE] (10pt, 30pt)--(10pt, 50pt); 
	\draw[thick, oldiagcolorE] (30pt, 50pt)--(50pt, 50pt); 
	\draw[thick, oldiagcolorE] (70pt, 50pt)--(70pt, 30pt); \begin{scope}[xshift=10pt, yshift=30pt]
	\draw [thick, rotate=270,black, l-system={Hilbert curve, step=20pt, angle=90, axiom=L, order=1}]  lindenmayer system;
	\end{scope}
	\begin{scope}[xshift=10pt, yshift=50pt]
	\draw [thick, black, l-system={Hilbert curve, step=20pt, angle=90, axiom=L, order=1}]  lindenmayer system;
	\end{scope}
	\begin{scope}[xshift=50pt, yshift=50pt]
	\draw [thick, black, l-system={Hilbert curve, step=20pt, angle=90, axiom=L, order=1}]  lindenmayer system;
	\end{scope}
	\begin{scope}[xshift=70pt, yshift=10pt]
	\draw [thick, rotate=90,black, l-system={Hilbert curve, step=20pt, angle=90, axiom=L, order=1}]  lindenmayer system;
	\end{scope}
	\end{tikzpicture}
\end{center}
As cores diferentes ajudam a explicar como $h_2$ foi construída. Colocamos primeiro cópias reduzidas do troço não-!!{colorC} de~$h_1$ no inferior esquerdo, superior esquerdo, superior direito e inferior direito do quadrado (traçadas a preto). Depois ligamos estas quatro figuras (com linhas !!{colorE}). Por fim, ligamos a figura ao contorno do quadrado (com linhas !!{colorC}).

Passamos a $h_3$. Tal como $h_2$ foi feita a partir de quatro cópias ligadas e reduzidas do troço não vermelho de~$h_1$, $h_3$ compõe-se de quatro cópias reduzidas do troço não vermelho de~$h_2$ (a preto), que depois se unem (com linhas !!{colorE}) e se ligam ao contorno do quadrado (com linhas !!{colorC}).
\begin{center}
	\begin{tikzpicture}
	\draw[gray] (0pt, 0pt) rectangle (80pt, 80pt);
	\draw[step = 10pt, gray, very thin] (0pt, 0pt) grid (80pt, 80pt);
	\draw[thick, oldiagcolorC](5pt, 0pt)--(5pt, 5pt); 
	\draw[thick, oldiagcolorC] (75pt, 0pt)--(75pt, 5pt); 
	\draw[thick, oldiagcolorE] (75pt, 45pt)--(75pt, 35pt);
	\draw[thick, oldiagcolorE] (5pt, 35pt)--(5pt, 45pt); 
	\draw[thick, oldiagcolorE] (35pt, 45pt)--(45pt, 45pt); 
	\draw[thick, oldiagcolorE] (75pt, 45pt)--(75pt, 35pt); \begin{scope}[xshift=5pt, yshift=35pt]
	\draw [rotate=270, thick, black, l-system={Hilbert curve, step=10pt, angle=90, axiom=L, order=2}]  lindenmayer system;
	\end{scope}
	\begin{scope}[xshift=5pt, yshift=45pt]
	\draw [thick, black, l-system={Hilbert curve, step=10pt, angle=90, axiom=L, order=2}]  lindenmayer system;
	\end{scope}
	\begin{scope}[xshift=45pt, yshift=45pt]
	\draw [thick, black, l-system={Hilbert curve, step=10pt, angle=90, axiom=L, order=2}]  lindenmayer system;
	\end{scope}
	\begin{scope}[xshift=75pt, yshift=5pt]
	\draw [thick, rotate=90,black, l-system={Hilbert curve, step=10pt, angle=90, axiom=L, order=2}]  lindenmayer system;
	\end{scope}
	\end{tikzpicture}
\end{center}
Agora vê-se o padrão geral para definir $h_{n+1}$ a partir de~$h_n$.
Por fim definimos a curva $h$ \emph{em si} considerando o
limite ponto a ponto destas funções sucessivas $h_1$, $h_2$,
\dots\@ Isto é, para cada $x \in \unitsquare$:
\[
	h(x) = \lim_{n \rightarrow \infty} h_n(x)
\]
Mostramos agora que esta curva preenche o espaço. Ao traçar $h_n$,
impomos uma grelha $2^n \times 2^n$ sobre~$\unitsquare$. Pelo teorema de Pitágoras, a diagonal de cada célula da grelha tem comprimento:
\[
\sqrt{\left(\nicefrac{1}{2^{n}}\right)^2+\left(\nicefrac{1}{2^{n}}\right)^2} = 2^{(\frac{1}{2}-n)}
\]
e é evidente que $h_n$ passa por todas as células. Logo cada ponto
de~$\unitsquare$ está \emph{a uma distância no máximo} $2^{(\frac{1}{2}-n)}$ de algum ponto de~$h_n$. Mas $h$ é definida como o limite das
funções $h_1$, $h_2$, $h_3$, \dots\@ Assim, a distância máxima de qualquer
ponto a~$h$ é dada por:
\[
\lim_{n \rightarrow \infty} 2^{(\frac{1}{2}-n)} = 0.
\]
Ou seja: todo o ponto de~$\unitsquare$ está à distância~$0$ de~$h$. Em
outras palavras, todo o ponto de~$\unitsquare$ situa-se \emph{sobre} a curva. Logo $h$
preenche o espaço!{}

Resta mostrar que $h$ é, de fato, uma \emph{curva}. Para isso
é preciso definir o conceito. A definição moderna baseia-se na de
Jordan, de 1887 (O.s.\@, poucos anos antes da primeira
curva preenchedora de espaço):

\begin{defn}
Uma curva é uma aplicação contínua de $\unitline$ para $\Real^2$.
\end{defn}

Isto é bastante intuitivo: uma curva é, intuitivamente, uma aplicação ``suave''
que leva uma reta canônica para o plano~$\Real^2$. A nossa função
$h$ é de fato uma aplicação de $\unitline$ para $\Real^2$. Logo só precisamos
de mostrar que $h$ é contínua. Definimos continuidade em
\olref[limits]{sec} com a notação $\epsilon$/$\delta$. Em linguagem corrente, queremos estabelecer o seguinte: \emph{Se especificar um
ponto $p$ em $\unitsquare$, juntamente com qualquer nível de
precisão $\epsilon$, podemos encontrar um intervalo aberto de $\unitline$ tal
que, dado qualquer $x$ nesse intervalo, $h(x)$ fica a uma distância inferior a~$\epsilon$
de~$p$.}

Assim: suponhamos que $p$ e~$\epsilon$ estão fixos. Isto corresponde, em
efeito, a traçar um círculo com centro $p$ e raio~$\epsilon$ sobre
$\unitsquare$. (O círculo pode sair para fora de~$\unitsquare$,
mas não importa.) Recordemos que, ao descrever
$h_n$, traçamos uma grelha $2^n \times 2^n$ sobre~$\unitsquare$. É
óbvio que, por pequeno que seja~$\epsilon$, existe algum~$n$
tal que alguma célula individual da grelha $2^n \times 2^n$
sobre~$\unitsquare$ fica inteiramente dentro do círculo com centro $p$ e
raio~$\epsilon$.

Tomemos esse~$n$, e seja $I$ a maior parte aberta de~$\unitline$
que $h_n$ mapeia inteiramente para a célula relevante. (É claro
que $I$ existe, pois já notamos que $h_n$ passa por
todas as células da grelha $2^n\times 2^n$.) Basta agora
mostrar que, sempre que $x \in I$, o ponto $h(x)$ pertence à mesma
célula. E, para \emph{isso}, basta mostrar que
$h_m(x)$ pertence à mesma célula, para qualquer $m > n$. Mas isto é
óbvio. Se considerarmos o que acontece com $h_m$ para $m > n$, vemos
que exatamente a ``mesma parte'' do intervalo unitário é mapeada
para a mesma célula; apenas a descrevemos de modo cada vez mais
esticado e sinuoso.

\end{document}
